% Part: computability
% Chapter: recursive-functions
% Section: halting-problem

\documentclass[../../../include/open-logic-section]{subfiles}

\begin{document}

\olfileid[fr]{cmp}{rec}{hlt}
\olsection{Le problème de l'arrêt}

Le \emph{problème de l'arrêt} consiste, sous sa forme générale,
à décider, étant donnés la description~$e$ d'une fonction
partielle calculable, par exemple son programme, et un entier~$n$,
si son calcul sur l'entrée~$n$ s'arrête, autrement dit s'il
produit un résultat. Alan Turing a montré que ce problème
ne peut lui-même être résolu par une fonction calculable :
la fonction
\[
h(e, n) =
\begin{cases}
1 & \text{si le calcul $e$ s'arrête sur l'entrée $n$}\\
0 & \text{sinon,}
\end{cases}
\]
n'est pas calculable.

Pour les fonctions récursives partielles, l'indice~$e$ du
théorème de forme normale de Kleene peut jouer le rôle
de description d'un programme. Si $f$ est une fonction
récursive partielle, tout entier~$e$ pour lequel l'équation
du théorème est satisfaite est un indice de~$f$. La
construction du théorème définit, pour \emph{tout}
entier~$e$, une fonction récursive partielle par
\[
\cfind{e}(x) \simeq U(\mu s \; T(e, x, s))
\]
Inversement, pour toute fonction récursive partielle unaire
$f\colon \Nat \to \Nat$, il existe un $e \in \Nat$ tel que
$\cfind{e}(x) \simeq f(x)$ pour tout $x \in \Nat$.
Il existe même une infinité de tels indices. La
\emph{fonction d'arrêt}~$h$ est définie par
\[
h(e, x) =
\begin{cases}
1 & \text{si $\cfind{e}(x) \fdefined$}\\
0 & \text{sinon.}
\end{cases}
\]
Ainsi, $h(e,x)=0$ exactement lorsque le calcul
$\cfind{e}(x)$ ne donne aucune valeur. Certains entiers
peuvent ne coder aucune description de programme dans un
codage particulier; ils désignent néanmoins ici une
fonction récursive partielle par la formule précédente.\footnote{Correction éditoriale : pour cette numérotation définie par $T$ et $U$, tout entier $e$ détermine $\cfind{e}$. Nous supprimons donc les alternatives où l'entier ne serait l'indice d'aucune fonction, présentes dans la source ; elles ne s'appliquent pas à cette définition.}

\begin{thm}
\ollabel{thm:halting-problem}
La fonction d'arrêt~$h$ n'est pas récursive partielle.
\end{thm}

\begin{proof}
Supposons que $h$ soit récursive partielle. Elle est totale
par définition; cette hypothèse en ferait donc une fonction
récursive. Définissons
\[
d(y) =
\begin{cases}
1 & \text{si $h(y, y) = 0$}\\
\umin{x}{x \neq x} & \text{sinon.}
\end{cases}
\]
Aucun entier~$x$ ne vérifie $x \neq x$; par conséquent,
$\umin{x}{x \neq x}$ est indéfini. Il s'ensuit que
\begin{enumerate}
\item $d(y) \fdefined$ si et seulement si
  $\cfind{y}(y) \fundefined$;
\item $d(y) \fundefined$ si et seulement si
  $\cfind{y}(y) \fdefined$.
\end{enumerate}
Comme $h$ est supposée récursive partielle et totale,
la définition par cas ci-dessus rendrait $d$ récursive
partielle. Pour le voir sans évaluer une branche indéfinie, on peut écrire $d(y)\simeq 1+\umin{z}{h(y,y)=0}$. Le théorème de forme normale lui donnerait
alors un indice~$e_d$. Examinons $h(e_d,e_d)$,
qui vaut nécessairement $0$ ou~$1$ :
\begin{enumerate}
\item Si $h(e_d,e_d)=1$, alors
  $\cfind{e_d}(e_d) \fdefined$. Mais
  $\cfind{e_d}\simeq d$, tandis que $d(e_d)$
  est défini si et seulement si $h(e_d,e_d)=0$ :
  contradiction.
\item Si $h(e_d,e_d)=0$, alors
  $\cfind{e_d}(e_d) \fundefined$. Or
  $\cfind{e_d}\simeq d$ et, par définition de~$d$,
  $d(e_d) \fdefined$ : contradiction.
\end{enumerate}
Les deux cas étant impossibles, $h$ ne peut être
récursive partielle.
\end{proof}

\end{document}
